%==============================================================================
% Capítulo 10: Conclusiones
%==============================================================================
\section{Conclusiones}\label{sec:conclusiones}

El trabajo desarrollado ha cumplido parcialmente el objetivo general de evaluar la viabilidad y el desempeño de la automatización del registro inicial de incidentes ---sin medir la intervención humana---, con el desarrollo y la implementación del sistema como medio: la combinación de N8N como orquestación, FastAPI como núcleo de procesamiento, PostgreSQL como persistencia, Twilio como canal telefónico y Gemini~3.6 Flash como motor de inferencia resultó técnicamente viable, eficiente en el registro y económicamente accesible para organizaciones medianas, con viabilidad operativa condicionada al balance neto (\autoref{subsec:balance-neto}). Ese cumplimiento no se extiende al segundo objetivo específico: la exactitud estricta del 73,5~\% y el macro-F1 de umbral (0,6508) quedaron por debajo de las cotas mínimas de 85~\% y 0,85 fijadas como criterio de éxito contingente, por lo que ese objetivo no se alcanza. La reducción del tiempo de registro del 71,2~\% sí se sostiene como resultado central, con significación estadística (véase la \autoref{sec:resultados}). Por sub-hipótesis, la hipótesis principal se confirma en tiempo de registro, se refuta en el umbral de clasificación y queda no contrastada en intervención humana por ausencia de registro primario; las dos hipótesis subsidiarias no se contrastaron (\autoref{subsec:interpretacion-resultados}).

\subsection{Cumplimiento de los objetivos del trabajo}\label{subsec:cumplimiento-objetivos}

La Tabla~\ref{tab:cumplimiento-objetivos} sintetiza el grado de cumplimiento de cada objetivo y la evidencia que lo respalda.

\begin{table}[htbp]
\centering
\caption{Cumplimiento de los objetivos del trabajo y evidencia de respaldo}
\label{tab:cumplimiento-objetivos}
\footnotesize
\setlength{\tabcolsep}{4pt}
\begin{tabular}{@{}>{\raggedright\arraybackslash}p{3.6cm} >{\raggedright\arraybackslash}p{2.2cm} >{\raggedright\arraybackslash}p{8.4cm}@{}}
\toprule
\textbf{Objetivo} & \textbf{Cumplimiento} & \textbf{Evidencia} \\
\midrule
Objetivo general: viabilidad y desempeño & Alcanzado parcialmente & Tiempo de registro alcanzado (71,2~\%); clasificación por debajo del umbral (objetivo 2); intervención humana no medida; viabilidad operativa condicionada (\autoref{subsec:balance-neto}) \\
Objetivo específico 1: arquitectura & Alcanzado parcialmente & \autoref{sec:arquitectura}; invocación sincrónica (\autoref{subsec:desafios-integracion}) y TLS no verificado (\autoref{subsec:vision-general}) \\
Objetivo específico 2: clasificador & No alcanzado & Exactitud 73,5~\%, IC [67,0~\%, 79,1~\%]; macro-F1 de umbral 0,6508 (0,5207 con cinco clases, una no estimable) \\
Objetivo específico 3: tres canales & Alcanzado parcialmente & \autoref{subsec:capa-canales}; latencia y exactitud por canal (Tabla~\ref{tab:exactitud-etapa-canal}); complejidad percibida no medida \\
Objetivo específico 4: comparación & Alcanzado parcialmente & \autoref{subsec:comparacion-tiempos}; exactitud solo del flujo automatizado (\autoref{subsec:balance-neto}); intervención humana no medida \\
Objetivo específico 5: documentación & Alcanzado & \autoref{sec:anexo-repositorio}; \autoref{sec:recomendaciones} \\
Objetivo específico 6: errores/calibración & Alcanzado & \autoref{subsec:analisis-errores}; \autoref{subsec:calibracion}; dentro de la muestra, \autoref{subsec:sensibilidad-ajustes} \\
\bottomrule
\end{tabular}
\par\smallskip
{\footnotesize\textit{Fuente:} elaboración propia.\par}
\end{table}

La arquitectura de cinco capas distribuyó las responsabilidades entre componentes comunicados mediante interfaces REST, y la elección de tecnologías maduras y autoalojadas (Docker, PostgreSQL, N8N) respondió a las restricciones presupuestarias, cambiarias y normativas de las organizaciones medianas argentinas. El primer objetivo específico se alcanza, sin embargo, solo parcialmente: la invocación de N8N al servicio de clasificación es sincrónica y la terminación TLS prevista no se verificó en la instancia publicada.

Sobre esa arquitectura se construyó el clasificador híbrido, cuyo desempeño sobre el corpus real quedó por debajo de las cotas mínimas fijadas en el segundo objetivo específico, por lo que ese objetivo no se alcanza; la brecha se analiza en la \autoref{sec:discusion}. La exactitud estricta, estimada dentro de la muestra, fue de 0,7405 (reglas) y 0,7246 (Gemini); la laxa, del 89,5~\%; por canal, de 0,654 a 0,811. La estrategia compositiva ---un filtro determinístico de primera etapa y un modelo de lenguaje grande reservado para los casos que el filtro no resuelve--- resultó el rasgo más relevante del sistema: el 65,5~\% de los casos se resuelve sin inferencia externa y el modelo actúa como árbitro de último recurso (\autoref{subsec:subsistema-clasificacion}); que esa estrategia sea más eficiente que un esquema solo de modelo no se contrastó (\autoref{subsec:interpretacion-resultados}).

Los tres canales quedaron integrados en un único flujo de orquestación. El canal telefónico, cuya transcripción automática del habla representaba el mayor desafío técnico, alcanzó un tiempo de registro medio de 20,1 segundos, que incluye la duración de la llamada grabada y es una cota inferior de la duración real por el tope de 45~s de grabación (\autoref{subsec:metricas-instrumentos}) ---un valor que lo posiciona como una vía prácticamente útil. La normalización de la entrada en el propio N8N, antes de invocar al servicio de clasificación, aisló la complejidad de cada protocolo de recepción del núcleo de procesamiento. Estos resultados son compatibles con la hipótesis subsidiaria de integración multicanal, pero no la contrastan: no se midieron la capacidad bajo carga concurrente ni la complejidad percibida, y el tercer objetivo específico se alcanza solo en su construcción.

La evaluación comparativa entre el flujo manual y el automatizado se presenta en la \autoref{subsec:comparacion-tiempos} y se interpreta, con sus alcances y límites, en la \autoref{subsec:interpretacion-resultados}. En síntesis, el sistema automatizado resultó más rápido que el operador humano en 188 de los 200 pares medidos, con la reducción del tiempo de registro como resultado central. La comparación no cubre la exactitud del brazo manual ni la intervención humana por caso, no contrastada por ausencia de registro primario, por lo que el cuarto objetivo específico se alcanza parcialmente y el aporte operativo asociado al gate de revisión se declara no medido: el sistema libera tiempo bruto de registro, cuyo saldo neto está indeterminado (\autoref{subsec:balance-neto}).

La documentación generada ---arquitectura, procedimientos de despliegue, especificación OpenAPI, flujo N8N exportable, análisis tipificado de errores y bandas de confianza--- constituye el aporte documental del quinto objetivo específico y sustenta el sexto. La tipología de tres patrones de error (descripciones cortas, mezcla de categorías y jerga local) ofrece un punto de partida concreto para replicar o adaptar el sistema y anticipar las categorías de fallo más probables.

\subsection{Aportes y generalización}\label{subsec:aportes-generalizacion}

Más allá del cumplimiento de los objetivos específicos, los resultados permiten formular una conclusión acotada a este corpus y al contraste efectivamente realizados: en este corpus de 200 incidentes de mesa de ayuda, en español rioplatense y con cinco categorías canónicas, un modelo de lenguaje grande invocado mediante instrucciones contextualizadas ---sin ajuste fino--- resultó una alternativa técnicamente viable para la clasificación de tickets cuando se combina con un filtro determinístico que absorba los casos triviales y reduzca el costo de inferencia. Esta lectura se limita a este dominio: no se realizó una comparación directa contra clasificadores supervisados entrenados sobre el mismo corpus, de modo que no se sostiene aquí una superioridad general de los modelos de lenguaje (véase la \autoref{subsec:limitaciones-estudio}).

Esta combinación arquitectónica ---orquestación visual, filtrado determinístico y LLM como árbitro semántico--- constituye un patrón de diseño reproducible para la gestión de servicios en organizaciones medianas de América Latina, donde la restricción presupuestaria y la soberanía de datos son determinantes. La arquitectura propuesta no depende de GPU ni de licencias propietarias; su costo monetario, dominado por la infraestructura y no por la inferencia, se cuantifica en la \autoref{subsec:implicancias-practicas}.

El trabajo aporta tres contribuciones que se sostienen más allá del caso particular. La primera es una arquitectura reproducible documentada en todos sus niveles ---modelo de datos, contrato REST y flujo de orquestación---, replicable por cualquier equipo con conocimientos básicos de Docker, Python y N8N. La segunda es evidencia empírica, bajo un diseño cuasiexperimental con muestras pareadas, doble etiquetado y prueba no paramétrica, sobre el desempeño de LLM en un dominio y una variante lingüística escasamente documentados; sus límites de independencia se declaran en la \autoref{subsec:limitaciones-estudio}. La tercera es un marco de implementación que contempla desde el inicio las restricciones normativas argentinas, sin requerir sortear barreras legales ni asumir riesgos de cumplimiento.

\newpage
